\section{Hermit Crabs and Nietzsche}

\paragraph{Perspectivism and the Human Race}
\emph{Part I of a series.}



Hermit crabs don't have a shell of their own. Since that makes them vulnerable to predators, they have to scramble to find empty shells left behind by gastropods. As they grow and no longer fit completely within their current shell, they once again are forced to go out and find a better fitting shell. 

\subparagraph{Ideation}
We find this also in the human being. Naked and vulnerable, with no real identity of his own, he goes out in search of a shell, which in this case, takes the form of an ideation: a conceptual scheme, system of thought, religion, philosophy, political theory, and so on. As his new ``home", it offers protection, shapes his identity and he defends it with vehemence. 

Such ideations are thus viewed from the perspective of each crab. As such, it seems perfectly reasonable and true to the crab. \textbf{Perspectivism}, as developed by Kant, is the philosophy that no such perspective is definitively true. Nietzsche developed this further and claimed each perspective is the manifestation of the ``\textbf{Will to Power}". 

So, rather than evaluating perspectives from their truth value, they should instead be evaluated by how much Power they manifest. Thus, a lower type of man, will adopt any old shell that he comes upon. Not thinking too much about it, he accepts his perspective as a given, usually on some sort of authority. 

A more ``thoughtful" person, hence with more power, is more selective. Having heard of higher ideals, this person looks for an ideational home that resembles those ideals. For this person, appearance trumps reality. Not fully understanding those ideals, he takes on the pose of what a spiritual person should look like, or a Christian, or a liberal, or a ``cool dude", and so on. 

Still higher types of men are the creators of ideational frameworks, and operate on a more serious level. These frameworks aspire to be rational, coherent, and evidential. Since few people are able to follow a logical argument, achieve consistency in thought, or make the effort to learn the facts, this approach appeals only to a few. 

\subparagraph{Predators}
A predator, thus, is anyone who is a threat to destroying one's ideational home. The third type will use the tools of debate to ward off a threat. So we see millions and millions of words defending one position and attacking others. The more skilled are able to take these arguments many levels deep, but ultimately, they end up in a stalemate. Some ideological systems, such as Marxism or Freudianism, are able to account for contradictory events, so no factual evidence is able to refute them, hence their appeal. 

The middle type is forced into a more emotional response. They will use ideals of ``love" or ``compassion" to counteract threats to their self-image. Usually, they promote a non-judgmentalism, which prevents threats from coming up in the first place. This group feels safe and secure among each other, and conflicts are reduced to the personal level, not ideational. Since coherence is not valued, different and even incompatible ideational schemes can be tolerated not only within the group, but even within a single individual. This gives the ability to ward off intellectual threats, but simply shifting one's own intellectual framework. 

\subparagraph{Resentment}
Threats to one's ideational ``home" or worldview give rise to resentment. If an ideational threat cannot be counteracted with a logical argument, or diffused with a claim to a higher spirituality, then resentment arises. This may take the forms of contempt, hostility, or even physical conflict. In some cases, one is forced to give up one home for another; this is called a ``\textbf{conversion}". Converts are valued in ideational battles. 

This happens despite the fact that they are all hermit crabs. There is a \textbf{forgetfulness of Being} or primary reality, which is replaced with the self-identification as a particular ``shell"; thus one lives in an artificial or \textbf{second reality} created by thought. 

\subparagraph{Man of Tradition}
The Traditional man lives free of illusion (not to be confused with cynicism or misanthropy). Although, as a being manifesting in the physical world, he necessarily has a perspective, he understands them for what they are. He remembers himself and knows who he is, transcendent to any ideational scheme. Instead of resentment, he has a sense of \emph{noblesse oblige}\footnote{\url{http://dictionary.reference.com/browse/noblesse-oblige}} towards others. Understanding the difficulties and efforts required to ``know oneself", he has compassion for those unable or struggling to achieve that state. However, he also understands he cannot make the efforts for them. 

\emph{Part II on the Will to Power is forthcoming.}



\flrightit{Posted on 2010-07-18 by Cologero }
